% Figure: normal-shock property ratios versus upstream Mach number M1,
% gamma=1.4. p2/p1, rho2/rho1, T2/T1 rise; M2 falls below 1.
\begin{figure}[htbp]
\centering
\begin{tikzpicture}
\begin{axis}[
  width=12cm, height=8cm,
  xlabel={upstream Mach number $M_1$}, ylabel={ratio across shock},
  xmin=1, xmax=4, ymin=0, ymax=20,
  axis lines=left,
  grid=major, grid style={enggray!20},
  tick label style={font=\scriptsize},
  label style={font=\small},
  legend style={font=\scriptsize, at={(0.03,0.97)}, anchor=north west},
  legend cell align=left,
]
% p2/p1 = 1 + (2g/(g+1))(M1^2 - 1), g=1.4 -> 1 + (2.8/2.4)(M^2-1) = 1 + 1.16667(M^2-1)
\addplot[engred, very thick, domain=1:4, samples=120] { 1 + 1.16667*(x^2 - 1) };
\addlegendentry{$p_2/p_1$}
% rho2/rho1 = (g+1)M^2 / ((g-1)M^2 + 2) = 2.4 M^2 / (0.4 M^2 + 2)
\addplot[engblue, very thick, domain=1:4, samples=120] { 2.4*x^2 / (0.4*x^2 + 2) };
\addlegendentry{$\rho_2/\rho_1$}
% T2/T1 = (p2/p1)*(rho1/rho2)
\addplot[engamber, very thick, domain=1:4, samples=120] {
  (1 + 1.16667*(x^2 - 1)) * (0.4*x^2 + 2) / (2.4*x^2) };
\addlegendentry{$T_2/T_1$}
% M2 = sqrt( ((g-1)M^2 + 2) / (2g M^2 - (g-1)) ) = sqrt((0.4M^2+2)/(2.8M^2-0.4))
\addplot[englime, very thick, domain=1:4, samples=120] {
  sqrt( (0.4*x^2 + 2) / (2.8*x^2 - 0.4) ) };
\addlegendentry{$M_2$ (downstream)}
% M2 = 1 reference
\draw[enggray, dashed] (axis cs:1,1) -- (axis cs:4,1);
\node[enggray, font=\scriptsize, anchor=south east] at (axis cs:4,1.0) {$M_2=1$};
% density limit annotation: rho2/rho1 -> 6 as M->inf
\draw[enggray, dotted] (axis cs:1,6) -- (axis cs:4,6);
\node[enggray, font=\scriptsize, anchor=south east] at (axis cs:4,6.0) {density limit $6$};
\end{axis}
\end{tikzpicture}
\caption[Normal-shock property ratios]{Property ratios across a normal
shock in air ($\gamma=1.4$) as functions of the upstream Mach number
$M_1$. Pressure (red) rises without bound. Temperature (amber) rises
without bound. Density (blue) rises but saturates at the limiting
ratio $(\gamma+1)/(\gamma-1)=6$ as $M_1\to\infty$: a normal shock
cannot compress a perfect gas beyond six times its upstream density,
which is why hypersonic shock layers stay thick relative to a naive
incompressible expectation. The downstream Mach number (green) falls
below $1$ for every $M_1>1$: the flow always leaves a normal shock
subsonic. The pressure, density, and temperature ratios are the
Rankine-Hugoniot relations; the downstream Mach number follows from
mass, momentum, and energy conservation across the discontinuity.}
\label{fig:vol03ch13:normal-shock}
\end{figure}
